\section{Lean 4: integers as equivalence classes of pairs}
\label{sec:lean-bridge-10}
Represent the formal difference $a-b$ by $(a,b)\in\mathbb N^2$.
Two pairs represent the same integer when $a+d=b+c$.
Reflexivity, symmetry, and transitivity are proved separately.
Addition of representatives is componentwise,
$(a,b)+(c,d)=(a+c,b+d)$; the last theorem shows that
equivalent inputs give equivalent outputs, making this operation
well-defined on equivalence classes.

\lstinputlisting[language=LeanFour,basicstyle=\ttfamily\scriptsize,
breaklines=true,columns=fullflexible,keepspaces=true,frame=single,
firstline=15,lastline=52]
{../../formal/lean/Con2physFormal/BridgesNumbers.lean}

\noindent\textit{Scope:} this code establishes an equivalence relation
and congruence of representative addition. It does not yet define
the quotient type or establish an explicit equivalence with \texttt{Int}.
